\clearpage
\VerdierPhysicalPageAnchor{0025}
ここで $\mathrm{RH}(X,{\ })$ は $X$ 上の「大域切断」関手の全右導来関手を
表し（$\mathrm{RH}(Y,{\ })$ と $\mathrm{RH}(X\times Y,{\ })$ も同様）、
$\overset{\mathrm L}{\otimes}_A$ はテンソル積関手の全左導来関手を表す。
（$X,Y$ は有限コホモロジー次元の空間、たとえば有限胞体複体であると
仮定する。）$\text{(Kü)}_{\mathrm{Der}}$ の両辺は、擬同型を除いて定まる
$A$-加群の複体であり、そのコホモロジー加群はそれぞれ
$\text{(Kü)}_{\mathrm{coh}}^{4}$ の第一および第二スペクトル系列の収束先
である。もちろん、$F,G$ が $X,Y$ 上の適切に有界な層の複体である場合も
$\text{(Kü)}_{\mathrm{Der}}$ は成立する。これは実際の計算で十分扱いやすく、
因子が複数あるときの可換性と結合性も定式化できる。

\VerdierSemanticAnchor{intro-5-base-change}
スカラー拡大は導来圏において類似の公式を与える。$A\to B$ を $A$-代数、
$F$ を $X$ 上の $A$-加群の層とすると、
\[
\mathrm{RH}(X,F)\overset{\mathrm L}{\otimes}_A B
\xrightarrow{\sim}
\mathrm{RH}\bigl(X,F\overset{\mathrm L}{\otimes}_A B\bigr)
\]
という関係がある。

\VerdierSemanticAnchor{intro-5-naive-formulas}
比喩的に言えば、一般に、\emph{伝統的には偽である素朴な公式も、導来圏で
扱えば真になる}。

\VerdierSemanticAnchor{intro-5-l-adic}
係数環 $\mathbb Z/\ell^r\mathbb Z$ の重要性は、とりわけスキームの
エタール・コホモロジーに由来する。この理論では本質的に捩れ層を扱わねば
ならないことが知られている [6]。環 $\mathbb Z/\ell^r\mathbb Z$ に値を
とるコホモロジーの逆極限をとることで、$\ell$-進整数環
$\widehat{\mathbb Z}_{\ell}$ に値をとる $\ell$-進コホモロジーが得られる。
公式 $\text{(Kü)}_{\mathrm{Der}}$ はこの状況でも成立し（たとえば $X,Y$
が代数閉体上の固有スキームで、$A=\mathbb Z/\ell^r\mathbb Z$ の場合）、
ただちに極限へ移って $\widehat{\mathbb Z}_{\ell}$-層について同様の公式
を与える。さらに $\widehat{\mathbb Z}_{\ell}$ は単項イデアル整域なので、
$\text{(Kü)}_{\mathrm{Der}}$ はコホモロジーをとると
$\text{(Kü)}_{\mathrm{coh}}^{3}$ 型の公式を与え、$\mathbb Q_{\ell}$ と
テンソルすると $\mathbb Q_{\ell}$-層について
$\text{(Kü)}_{\mathrm{coh}}^{1}$ 型の公式を与える。

\VerdierSemanticAnchor{intro-6}
\noindent\textbf{6.}\quad
本論文では、以上に概説した考察を体系的に展開する。その構成は次のとおり
である。

\VerdierSemanticAnchor{intro-6-plan-chapter-one}
第 I 章では加法圏の複体の圏を研究する。これらの圏は次数付けられている。
より正確には、シフト関手（次数を所定の方向へ移し、それに応じて微分の
符号を変える）を備えている。
